\chapter{Может ли суперсвязность вывести Skyrme-коэффициент}

Проверяются три известных механизма появления четырёхпроизводного члена.

Для формы Маурера--Картана $L=V^{-1}dV$ выполнено
\begin{equation}
 dL+L\wedge L=0.
\end{equation}
Поэтому обычный Yang--Mills квадрат её кривизны тождественно нулевой и не
может породить недостающий $E_4$.

Для проектора $P$ каноническая грассманова связность
\begin{equation}
 A_P=P\,dP,
 \qquad F_P=P(dP)^2
\end{equation}
имеет ненулевую кривизну. Квадрат $\Tr(F_P^*F_P)$ даёт правильный
четырёхпроизводный масштаб $R^{-1}$. Но он требует пространственного
хопфова проектора и является отдельным инвариантом от уже выведенного
$\tau([N,V]^*[N,V])$. Общий след нормирует каждый после выбора метрики,
но не фиксирует размерное отношение между ними.

Третий известный путь --- редукция Yang--Mills теории в дополнительном
измерении. В литературе Skyrme-член действительно возникает вместе с
квадратичным sigma-model членом, но коэффициенты являются интегралами от
профиля дополнительной координаты и сопровождаются башней векторных мод;
см. \path{arXiv:2209.06607v2} и \path{arXiv:hep-th/0508130}.
В текущей версии такой профиль и дополнительное измерение отсутствуют.

Суперсвязность $\mathbb A=\nabla+\Phi$ упаковывает $F$, $D\Phi$ и
$\Phi^2$, но масштабирование $\Phi\mapsto a\Phi$ меняет их относительные
веса. Это точно совпадает с ранее доказанным пределом следа $M_{35}$:
после внешнего задания пространственной геометрии действие допустимо, но
сам след не выводит её масштаб и вложение.

\begin{theorem}[Разрыв Skyrme-коэффициента]
Текущая суперсвязность не выводит одновременно $E_2$ и $E_4$ с
фиксированным относительным коэффициентом. Проекторная кривизна является
математически правильным кандидатом для $E_4$, но её пространственный
носитель и общий масштаб с петлевым действием остаются дополнительными
данными.
\end{theorem}

\begin{protocol}
\begin{enumerate}
 \item чистая форма Маурера--Картана даёт $E_4$: \textbf{нет};
 \item проекторная кривизна ненулевая: \textbf{да};
 \item её коэффициент связан с $E_2$ текущим следом: \textbf{нет};
 \item высшеразмерный Yang--Mills выводит оба члена: \textbf{условно да};
 \item дополнительный профиль присутствует в проекте: \textbf{нет};
 \item конечный радиус версии V: \textbf{не выведен};
 \item следующий гейт: \textbf{общий масштаб проекторной суперсвязности}.
\end{enumerate}
\end{protocol}

Машинный сертификат сохранён в
\path{s2t_v5_superconnection_skyrme_coefficient_gate_results.json}.